\documentclass[12pt]{article}

\title{Multi-Page Sample with Table}
\author{PyTeX Workspace}
\date{\today}

\begin{document}
\maketitle

\section*{Introduction}
This document is a simple multi-page LaTeX sample.
It includes regular prose and a table float.
The explicit page breaks ensure the output spans multiple pages.

A workflow often begins with rough planning text.
This paragraph is intentionally plain and repetitive so the sample stays
portable and easy to compile in minimal LaTeX setups.

When drafting reports, teams often sketch assumptions, list constraints,
and write short status notes before converting everything into polished
technical sections.

\newpage

\section*{Data Snapshot}
The table below summarizes example milestones across three quarters.

\begin{table}[htbp]
\centering
\caption{Quarterly Milestones}
\begin{tabular}{|l|c|c|c|}
\hline
Item & Q1 & Q2 & Q3 \\
\hline
Prototype builds & 3 & 5 & 7 \\
Integration tests & 12 & 18 & 26 \\
Bug fixes merged & 9 & 14 & 21 \\
Documentation updates & 4 & 8 & 11 \\
\hline
\end{tabular}
\end{table}

The remaining text keeps this document realistic as a page-layout test.
It demonstrates that the same source can combine narrative sections and
floating material without any special packages.

Additional body text is included here to represent commentary around the
numbers in the table: trends, risks, and planned next actions.

\newpage

\section*{Conclusion}
This final page confirms that the file spans multiple pages.
It also serves as a minimal regression input for table placement behavior.

If needed, you can extend this sample by adding more rows to the table,
more sections, or bibliography entries.

\end{document}
